\section*{Bab \ref{ch:b2:diffeq} --- Persamaan Diferensial}

\begin{solution}{exo:b2:diffeq:1}
Matriks pertama: \omterm{def:b2:reduction:eigen}{nilai eigennya} $1, 2$, dengan \omterm{def:b2:reduction:eigen}{vektor eigen} $(1,0)$ dan $(1,1)$.
Uraikan $X_0 = (1,0) = 1\cdot(1,0) + 0\cdot(1,1)$: jadi penyelesaiannya
adalah
\[
X(t) = \eu^{t}\begin{pmatrix}1\\ 0\end{pmatrix} .
\]
(Sebab vektor awalnya sendiri sebuah \omterm{def:b2:reduction:eigen}{vektor eigen}.)

Yang kedua: $A = 2I + N$ dengan $N = E_{12}$ dan $N^2 = 0$: jadi $\eu^{tA} =
\eu^{2t}(I + tN)$, sehingga
\[
X(t) = \eu^{2t}\begin{pmatrix} 1 & t\\ 0 & 1\end{pmatrix}
\begin{pmatrix}1\\ 0\end{pmatrix}
= \eu^{2t}\begin{pmatrix}1\\ 0\end{pmatrix}.
\]
\end{solution}

\begin{solution}{exo:b2:diffeq:2}
$y' = \sin(ty)$: sebab $\abs{\sin(ty) - \sin(tz)} \leq \abs t\,\abs{y -
z}$ --- jadi \omterm{def:b2:metric:continuity}{Lipschitz} dalam $y$ secara \omterm{def:b2:funcseq:def}{seragam} pada setiap ruas waktunya:
sehingga penyelesaian global yang tunggal pada $\R$ (terapkanlah teoremanya pada setiap
ruasnya).

$y' = y^2$: hanya \omterm{def:b2:metric:continuity}{Lipschitz} secara lokal; jadi tak ada teorema global, dan memang
$y(0) = 1$ meledak di $t = 1$.

$y' = \abs y$: sebab $\abs\cdot$ bersifat \omterm{def:b2:metric:continuity}{Lipschitz} berkonstanta $1$: jadi keberadaan dan
ketunggalan global. Dengan $y(0) = 0$: $y \equiv 0$ (berkat ketunggalannya!). Dengan $y(0)
= 1$: $y$ tetap positif (sebab ia tak dapat menyeberangi penyelesaian nolnya), jadi $y' =
y$: sehingga $y = \eu^t$.
\end{solution}

\begin{solution}{exo:b2:diffeq:3}
Bila $y(t_1) = z(t_1)$ pada suatu waktu, maka $y$ dan $z$ menyelesaikan masalah
Cauchy yang sama di $t_1$: jadi ketunggalannya memaksa $y = z$ pada interval
bersamanya --- sehingga penyelesaian yang berbeda tak pernah bertemu.

Untuk $y' = y(1-y)$: konstanta $0$ dan $1$ merupakan penyelesaiannya. Sebuah
penyelesaian yang bermula di $\intoo{0}{1}$ tak pernah dapat mencapai $0$ atau $1$
(sebab ia akan bertabrakan dengan penyelesaian tetapnya): jadi ia tinggal di
$\intoo{0}{1}$, sehingga bersifat global (sebab terbatas, jadi tanpa ledakan --- misalnya
lewat kriteria \cref{exo:b2:diffeq:9}, atau karena medan vektornya
terbatas pada jalur yang menjebaknya).
\end{solution}

\begin{solution}{exo:b2:diffeq:4}
$(A - 2I)^2 = \begin{pmatrix}1 & 1\\ -1 & -1\end{pmatrix}^2 = 0$: jadi
Dunford dengan $D = 2I$ dan $N = A - 2I$:
\[
\eu^{tA} = \eu^{2t}\,(I + tN)
= \eu^{2t}\begin{pmatrix} 1 + t & t\\ -t & 1 - t\end{pmatrix}.
\]
Variasi konstanta dengan $B(s) = (\eu^{2s}, 0)^{\mathsf T}$:
\[
X(t) = \int_0^t \eu^{(t-s)A}B(s)\,\dd s
= \eu^{2t}\int_0^t \begin{pmatrix} 1 + (t-s)\\ -(t-s)
\end{pmatrix}\dd s
= \eu^{2t}\begin{pmatrix} t + \frac{t^2}{2}\\[2pt] -\frac{t^2}{2}
\end{pmatrix},
\]
dengan memakai $\eu^{(t-s)A}B(s) = \eu^{2(t-s)}(I + (t-s)N)\,(\eu^{2s},
0)^{\mathsf T} = \eu^{2t}\bigl(1 + (t-s),\, -(t-s)\bigr)^{\mathsf
T}$. (Periksa: $X(0) = 0$; dan $X' - AX = (\eu^{2t}, 0)^{\mathsf T}$
lewat penurunan.)
\end{solution}

\begin{solution}{exo:b2:diffeq:5}
\emph{Lewat deretnya:} $A = \alpha I + \beta J$ dengan $J =
\begin{pmatrix}0 & -1\\ 1 & 0\end{pmatrix}$ dan $J^2 = -I$; kedua
sukunya komut, jadi $\eu^{tA} = \eu^{\alpha t}\,\eu^{\beta tJ}$,
lalu deret $\eu^{\beta t J}$ terpecah menurut pangkat genap dan ganjilnya
menjadi $\cos(\beta t)I + \sin(\beta t)J$: yakni matriks perputaran-penskalaan
yang dinyatakan.

\emph{Lewat kompleksnya:} kenalilah $(x, y) \in \R^2$ dengan $z = x + \iu y$;
maka sistem $X' = AX$ berbunyi $z' = (\alpha + \iu\beta)z$, yang
penyelesaiannya $z(t) = \eu^{\alpha t}\eu^{\iu\beta t}z_0$ persis merupakan
spiralnya: dengan modulus $\eu^{\alpha t}$ dan argumen yang maju berkelajuan
$\beta$.
\end{solution}

\begin{solution}{exo:b2:diffeq:6}
Misalkan $v(t) = \eu^{-k(t - t_0)}\int_{t_0}^t u$. Maka
\[
v'(t) = \eu^{-k(t-t_0)}\Bigl(u(t) - k\int_{t_0}^t u\Bigr)
\leq C\,\eu^{-k(t-t_0)},
\]
menurut hipotesisnya. Mengintegralkannya dari $t_0$ ke $t$ (dengan $v(t_0) = 0$):
$v(t) \leq \frac{C}{k}\bigl(1 - \eu^{-k(t - t_0)}\bigr)$, yakni
$k\int_{t_0}^t u \leq C\bigl(\eu^{k(t-t_0)} - 1\bigr)$; lalu menyuapkan
ini kembali ke hipotesisnya: $u(t) \leq C\eu^{k(t-t_0)}$.

Ketunggalan dan kebergantungannya: dua penyelesaian $y, z$ persamaan
integralnya memenuhi
\[
\norm{y(t) - z(t)} \leq \norm{y_0 - z_0} +
k\int_{t_0}^{t}\norm{y - z},
\]
dan Gronwall dengan $C = \norm{y_0 - z_0}$ memberi batas
eksponensialnya; sedangkan $C = 0$ memberi ketunggalannya.
\end{solution}

\begin{solution}{exo:b2:diffeq:7}
Sulihkan $y = y_1 z$ dengan $y_1 = \frac{\sin t}{t}$: maka rumus
penurunan umumnya (\cref{prop:b2:diffeq:secondorder}) memberi, untuk $u
= z'$,
\[
y_1 u' + \Bigl(2y_1' + \frac{2}{t}\,y_1\Bigr)u = 0
\]
(dengan persamaannya dinormalkan menjadi $y'' + \frac2t y' + y = 0$). Hitunglah
$2y_1' + \frac2t y_1 = 2\,\frac{t\cos t - \sin t}{t^2} +
\frac{2\sin t}{t^2} = \frac{2\cos t}{t}$: jadi
\[
\frac{u'}{u} = -\frac{2\cos t}{t}\cdot\frac{t}{\sin t}
= -2\cot t
\quad\Longrightarrow\quad
u = \frac{1}{\sin^2 t}
\quad (\text{hingga sebuah konstanta}),
\]
lalu $z = -\cot t$, sehingga memberi $y_2 = y_1 z = -\frac{\cos t}{t}$.
Penyelesaian umumnya pada $\intoo{0}{\pi}$:
\[
y(t) = \alpha\,\frac{\sin t}{t} + \beta\,\frac{\cos t}{t} .
\]
(Fungsi ini adalah fungsi Bessel sferis berorde nol.)
\end{solution}

\begin{solution}{exo:b2:diffeq:8}
Dunford: $A = D + N$ yang komut, dengan $D$ \omterm{def:b2:reduction:diag}{dapat didiagonalkan} bernilai eigen
sama, dan $N$ nilpoten, jadi
\[
\eu^{tA} = \eu^{tD}\,\eu^{tN},
\qquad
\eu^{tN} = \sum_{k < n} \frac{t^kN^k}{k!}
\ \text{(sebuah polinomial matriks dalam } t).
\]
Misalkan $-2\alpha = \max_i \Re\lambda_i < 0$. Dalam basis yang mendiagonalkan
$D$, berlaku $\vertiii{\eu^{tD}} \leq \eu^{-2\alpha t}$ (sebab entrinya
$\eu^{t\lambda_i}$ bermodulus $\eu^{t\Re\lambda_i}$); sedangkan \omterm{def:b2:nvs:norm}{norma} pada
basis yang berbeda hanya berselisih konstanta. Karena itu
\[
\norm{X(t)} \leq \vertiii{\eu^{tA}}\,\norm{X_0}
\leq C'\,\eu^{-2\alpha t}\,(1 + t)^{n-1}\,\norm{X_0}
\leq C\,\eu^{-\alpha t}\norm{X_0} ,
\]
dengan menyerap polinomialnya ke dalam satu faktor eksponensial
(sebab $\eu^{-\alpha t}(1+t)^{n-1} \to 0$, sehingga terbatas).
\end{solution}

\begin{solution}{exo:b2:diffeq:9}
Misalkan $y$ penyelesaian maksimal pada $\intco{0}{T}$ dengan $T \leq \infty$,
lalu andaikan $T < \infty$. Bentuk integralnya memberi, untuk $t < T$,
\[
\norm{y(t)} \leq \norm{y_0} + \int_0^t \bigl(a\norm{y(s)} +
b\bigr)\dd s
\leq \bigl(\norm{y_0} + bT\bigr) + a\int_0^t\norm{y(s)}\,\dd s ,
\]
dan Gronwall membatasi $\norm{y(t)} \leq (\norm{y_0} +
bT)\,\eu^{aT} =: M$ pada $\intco{0}{T}$: jadi penyelesaiannya tinggal di sebuah
bola \omterm{def:b2:metric:compact}{kompak}. Maka $y' = f(t, y)$ terbatas di dekat $T$, jadi $y$ bersifat
\omterm{def:b2:metric:continuity}{Lipschitz} di dekat $T$ sehingga meluas secara \omterm{def:b2:metric:continuity}{kontinu} ke $T$ (lewat kriteria
Cauchy); lalu menyelesaikan masalah Cauchy di $(T, y(T))$ memperpanjang $y$
melampaui $T$, yang bertentangan dengan kemaksimalannya. Karena itu $T = \infty$: jadi tak ada
pelarian dalam waktu berhingga di bawah pertumbuhan linear.
\end{solution}

\begin{solution}{exo:b2:diffeq:10}
Akar karakteristik $r^2 - 3r + 2$: yakni $1$ dan $2$, jadi
penyelesaian homogennya adalah $a\eu^t + b\eu^{2t}$. Adapun tebakan
$\alpha\eu^t$ gagal karena $\eu^t$ sudah menyelesaikan
persamaan homogennya (sebab akar $r = 1$ ``beresonansi'' dengan
ruas kanannya). Dengan $y = \alpha t\,\eu^t$: $y' = \alpha(1 +
t)\eu^t$, $y'' = \alpha(2 + t)\eu^t$, dan
\[
y'' - 3y' + 2y = \alpha\eu^t\bigl(2 + t - 3 - 3t + 2t\bigr)
= -\alpha\,\eu^t :
\]
jadi $\alpha = -1$ dan $y_p = -t\,\eu^t$. Penyelesaian umumnya: $y = a\eu^t
+ b\eu^{2t} - t\eu^t$. Dengan data awal $y(0) = y'(0) = 0$: $a + b
= 0$ dan $a + 2b - 1 = 0$, jadi $b = 1$ dan $a = -1$:
\[
y(t) = \eu^{2t} - (1 + t)\,\eu^{t} .
\]

\end{solution}

\begin{solution}{exo:b2:diffeq:11}
$A = \lambda I + N$ dengan $N = E_{12} + E_{23}$: jadi $N^2 = E_{13}$,
$N^3 = 0$, dan $\lambda I$ komut dengan $N$:
\[
\eu^{tA} = \eu^{\lambda t}\Bigl(I + tN +
\frac{t^2}{2}N^2\Bigr)
= \eu^{\lambda t}\begin{pmatrix}
1 & t & \frac{t^2}{2}\\
0 & 1 & t\\
0 & 0 & 1
\end{pmatrix}.
\]
Penyelesaiannya: $X(t) = \eu^{\lambda t}\bigl(X_0 + tNX_0 +
\frac{t^2}2N^2X_0\bigr)$ --- jadi tiap komponennya adalah $\eu^{\lambda
t}$ kali polinomial berderajat $\leq 2$. Adapun batas derajatnya adalah
indeks kenilpotenannya dikurangi satu: sebab deret $\eu^{tN}$
terpotong di $N^2$.
\end{solution}

\begin{solution}{exo:b2:diffeq:12}
\begin{enumerate}
  \item Bila $X(T) = X(0)$, maka $Y(t) = X(t + T)$ menyelesaikan $Y' =
        AY + B(t + T) = AY + B(t)$ dengan $Y(0) = X(0)$: jadi
        ketunggalannya (\cref{thm:b2:diffeq:linear}) memberi $Y = X$,
        yakni $X$ berkala-$T$. Sedangkan konversnya trivial.
  \item Trigonalkan atas $\C$: $A = PT'P^{-1}$ dengan $T'$ segitiga
        atas dan berdiagonal $(\lambda_i)$. Setiap pangkat matriks
        segitiga bersifat segitiga dengan diagonal
        $(\lambda_i^k)$, jadi $\eu^{TA} = P\eu^{TT'}P^{-1}$ bersifat
        segitiga dalam basis yang sama dengan diagonal
        $(\eu^{T\lambda_i})$: itulah \omterm{def:b2:reduction:eigen}{nilai eigennya}. Lalu
        $I - \eu^{TA}$ terbalikkan bila dan hanya bila $\eu^{T\lambda} \neq
        1$ untuk setiap \omterm{def:b2:reduction:eigen}{nilai eigennya}, yakni bila $T\lambda \notin
        2\iu\pi\Z$, dan itulah hipotesis yang dinyatakan.
  \item Variasi konstanta: $X(T) = \eu^{TA}X(0) +
        \int_0^T\eu^{(T-s)A}B(s)\dd s$, jadi $X(T) = X(0)$ berbunyi
        \[
        \bigl(I - \eu^{TA}\bigr)X(0) =
        \int_0^{T}\eu^{(T-s)A}B(s)\,\dd s ,
        \]
        yang punya penyelesaian tunggal $X(0)$ di bawah hipotesis
        keterbalikannya: jadi tepat satu penyelesaian berkala-$T$.
        Adapun untuk osilator harmoniknya (dengan $\lambda =
        \pm\iu\omega$), kasus yang dikecualikan adalah $\omega T \in
        2\pi\Z$: yakni pemaksaan yang kalanya kelipatan
        kala alaminya --- itulah \omterm{pb:b2:diffeq:1}{resonansi}, sebagaimana dikuantifikasi
        soal akhir pekan.
\end{enumerate}
\end{solution}

\begin{solution}{pb:b2:diffeq:1}
\textbf{1.} \omterm{def:b2:reduction:charpoly}{Polinomial karakteristiknya} adalah $\lambda^2 -
\tau\lambda + \delta$, dengan akar
$\frac{\tau\pm\sqrt\Delta}{2}$. Bila $\delta < 0$ maka $\Delta =
\tau^2 - 4\delta > 0$ dan kedua akar realnya berhasil kali
$\delta < 0$: jadi bertanda berlawanan. Bila $\delta > 0$ dan $\Delta \geq
0$: akar realnya berhasil kali $> 0$ dan berjumlah $\tau$: jadi keduanya bertanda
sama dengan $\tau$. Bila $\Delta < 0$: sepasang sekawan $\alpha \pm
\iu\beta$ dengan $\alpha = \frac\tau2$ dan $\beta =
\frac{\sqrt{-\Delta}}2$.

\textbf{2.} $X(t) = a\,\eu^{\mu t}v_- + b\,\eu^{\lambda t}v_+$.
Orbit dengan $a = 0$ (masing-masing $b = 0$) berjalan sepanjang garis eigen
takstabilnya (masing-masing yang stabil); sedangkan yang lain punya $\norm X \to \infty$
pada kedua arah waktunya, dan asimtotik terhadap $\R v_+$ ketika $t \to
+\infty$ serta terhadap $\R v_-$ ketika $t \to -\infty$: itulah gambaran
pelananya. Adapun keterbatasan pada seluruh $\R$ memaksa $b = 0$ (sebab bila tidak, terjadi
ledakan di $+\infty$) dan $a = 0$ (di $-\infty$): jadi hanya
titik asalnya.

\textbf{3.} Dengan $\mu < \lambda < 0$, kedua eksponensialnya meluruh, jadi
$X(t) \to 0$. Bila $b \neq 0$, faktorkanlah $\eu^{\lambda t}$:
\[
X(t) = \eu^{\lambda t}\bigl(b\,v_\lambda + a\,\eu^{(\mu -
\lambda)t}v_\mu\bigr),
\qquad \eu^{(\mu-\lambda)t} \to 0 :
\]
jadi arah $X(t)$ menuju $\R v_\lambda$, yakni arah eigen yang
\emph{lambat} --- sehingga semua orbit selain sumbu cepatnya
tiba secara menyinggungnya (lihatlah panel kanan pada potret fasa
bab ini).

\textbf{4.} Dalam basis tempat $A = \begin{pmatrix} \alpha &
-\beta\\ \beta & \alpha\end{pmatrix}$
(\cref{exo:b2:diffeq:5}; dan untuk sistem osilator Bagian II
bentuk ini dicapai lewat penggantian basis real yang gamblang),
penyelesaiannya adalah $\eu^{\alpha t}$ kali perputaran bersudut $\beta
t$: yakni spiral logaritmik, yang mengerut ketika $\alpha = \frac\tau2
< 0$, memuai ketika $\tau > 0$, dan berupa kurva tertutup (yakni elips
dalam koordinat aslinya) ketika $\tau = 0$: itulah pusatnya.

\textbf{5.} Syarat $\Delta = 0$ memberi \omterm{def:b2:reduction:eigen}{nilai eigen} ganda $\lambda =
\frac\tau2$; jadi menurut Cayley--Hamilton
(\cref{thm:b2:reduction:cayleyhamilton}), $(A - \lambda I)^2 =
0$, sehingga $N = A - \lambda I$ bersifat nilpoten, komut dengan
$\lambda I$, dan $\eu^{tA} = \eu^{\lambda t}(I + tN)$. Bila $N =
0$: maka $A = \lambda I$ dan semua sinarnya berupa orbit (yakni simpul bintang). Bila $N
\neq 0$: $X(t) = \eu^{\lambda t}(X_0 + tNX_0)$, dan untuk $NX_0
\neq 0$ arahnya konvergen ke satu-satunya arah eigennya
$\operatorname{im}N$: yakni simpul takwajar. Dan ini merampungkan
gambaran trace--determinannya.

\textbf{6.} Berlaku $\tau = -2\zeta\omega$, $\delta = \omega^2 > 0$,
$\Delta = 4\omega^2(\zeta^2 - 1)$. Jadi: $0 < \zeta < 1$ memberi
$\Delta < 0$ dan $\tau < 0$, yakni spiral stabil; $\zeta = 1$: $\Delta =
0$, yakni simpul stabil yang merosot; $\zeta > 1$: $\Delta > 0$, $\tau <
0$, $\delta > 0$, yakni simpul stabil; dan $\zeta = 0$: $\tau = 0$,
$\delta > 0$, yakni pusat. Jadi sebuah perjalanan tegak pada bidangnya di
$\delta = \omega^2$.

\textbf{7.} Akarnya $r = -\zeta\omega \pm
\omega\sqrt{\zeta^2-1}$. Untuk $\zeta < 1$: $r = -\zeta\omega
\pm \iu\omega_d$ dengan $\omega_d = \omega\sqrt{1-\zeta^2}$:
\[
x(t) = \eu^{-\zeta\omega t}\bigl(a\cos\omega_dt +
b\sin\omega_dt\bigr)
= R\,\eu^{-\zeta\omega t}\cos(\omega_dt - \varphi) .
\]
Untuk $\zeta = 1$: $x = (a + bt)\eu^{-\omega t}$. Untuk $\zeta >
1$: $x = a\eu^{r_-t} + b\eu^{r_+t}$, dengan kedua lajunya negatif.
Adapun maksimum berurutan $\abs x$ pada kasus teredam kurangnya terjadi pada
waktu yang terpisah sejauh pseudokalanya $\frac{2\pi}{\omega_d}$
(sebab fasa kosinusnya sama), dan nisbahnya adalah
$\eu^{-\zeta\omega\cdot2\pi/\omega_d} =
\eu^{-2\pi\zeta/\sqrt{1-\zeta^2}}$: yakni dekremen logaritmiknya,
sebuah alat ukur peredaman yang terbaca pada osiloskop.

\textbf{8.} Setelah dirasionalkan,
\[
\abs{\lambda_{\mathrm{lambat}}} = \omega\bigl(\zeta -
\sqrt{\zeta^2-1}\bigr)
= \frac{\omega}{\zeta + \sqrt{\zeta^2 - 1}} ,
\]
yang penyebutnya naik terhadap $\zeta \geq 1$: jadi laju
peluruhannya terbesar di $\zeta = 1$, tempat ia sama dengan $\omega$. Sebuah
pintu yang teredam lebih menutup tanpa membanting tetapi \emph{lambat};
jadi peredaman kritis adalah optimum bagi sang insinyur.

\textbf{9.} Berlaku $E' = x'x'' + \omega^2xx' =
x'\bigl(-2\zeta\omega x' - \omega^2x\bigr) + \omega^2xx' =
-2\zeta\omega\,x'^2 \leq 0$. Andaikan $x$ berkala dan
tak tetap, maka $E$ akan berkala sekaligus tak naik, sehingga
tetap, yang memaksa $x' \equiv 0$: jadi $x$ tetap, lalu
$\omega^2x = 0$: sehingga $x \equiv 0$. Jadi untuk $\zeta > 0$ satu-satunya
penyelesaian berkalanya adalah istirahat: sebab peredaman membunuh setiap daurnya.

\textbf{10.} Pusat $\zeta = 0$ hidup pada garis $\tau =
0$ pada bidang trace--determinannya --- yakni himpunan yang interiornya
kosong: jadi usikan matriksnya yang sekecil apa pun
(yakni peredaman fisis apa pun) menggeser $\tau$ dari nol lalu mengubah
orbit tertutupnya menjadi spiral. Karena itu kekalaan osilator
tak teredamnya merupakan gejala di mata pisau, bukan yang tegar.

\textbf{11.} Setelah menyulihkan $x_p = \Re(z\eu^{\iu\gamma t})$ ke dalam
persamaannya:
\[
\bigl(-\gamma^2 + 2\iu\zeta\omega\gamma +
\omega^2\bigr)z = F
\quad\Longrightarrow\quad
z = \frac{F}{\omega^2 - \gamma^2 + 2\iu\zeta\omega\gamma},
\]
jadi $x_p = \abs z\cos(\gamma t - \varphi)$ dengan $\varphi =
\arg(\omega^2 - \gamma^2 + 2\iu\zeta\omega\gamma)$, yakni
$\tan\varphi = \frac{2\zeta\omega\gamma}{\omega^2 -
\gamma^2}$, beserta $A(\gamma) = \abs z$ yang dinyatakan.

\textbf{12.} Selisih dua penyelesaiannya menyelesaikan persamaan
homogennya, yang untuk $\zeta > 0$ meluruh ke $0$
(pertanyaan 7): jadi setiap penyelesaiannya sama dengan $x_p$ ditambah transien
yang lenyap di tak hingga. Jadi keadaan tunaknya menjadi penarik global:
sebab syarat awalnya terlupakan, dan hanya $A(\gamma)$ beserta
ketertinggalan fasa $\varphi$ yang tersisa.

\textbf{13.} Minimumkan $g(u) = (\omega^2 - u)^2 +
4\zeta^2\omega^2u$ atas $u = \gamma^2 \geq 0$: maka $g'(u) =
-2(\omega^2 - u) + 4\zeta^2\omega^2 = 0$ di $u = \omega^2(1 -
2\zeta^2)$, yang di dalam bila dan hanya bila $\zeta < \frac{1}{\sqrt2}$. Di sana
\[
g(u_*) = 4\zeta^4\omega^4 + 4\zeta^2\omega^4(1 - 2\zeta^2)
= 4\zeta^2\omega^4(1 - \zeta^2),
\qquad
A(\gamma_*) = \frac{F}{2\zeta\omega^2\sqrt{1 - \zeta^2}} .
\]
Terhadap tanggapan statisnya $A(0) = \frac{F}{\omega^2}$:
penguatannya $\frac{1}{2\zeta\sqrt{1-\zeta^2}} \approx
\frac{1}{2\zeta}$ untuk $\zeta$ yang kecil --- jadi sistem yang teredam ringan
di dekat $\gamma_* \approx \omega$ melipatgandakan masukannya
seratus kali ketika $\zeta = 0.005$.

\textbf{14.} Fungsi $x$ yang dinyatakan memenuhi $x(0) = x'(0) = 0$ dan
\[
x'' + \omega^2 x = \frac{F(\omega^2 -
\gamma^2)\cos\gamma t}{\omega^2 - \gamma^2} = F\cos\gamma t
\]
(sebab bagian $\cos\omega t$-nya saling meniadakan). Adapun bentuk hasil kalinya menyusul
dari $\cos p - \cos q = 2\sin\frac{q+p}{2}\sin\frac{q-p}2$
dengan $p = \gamma t$ dan $q = \omega t$. Untuk $\gamma$ yang dekat ke
$\omega$, faktor $\sin\frac{(\omega-\gamma)t}2$ merupakan selubung
lambat yang memodulasi ayunan cepat
$\sin\frac{(\omega+\gamma)t}2$: yakni layangan, dengan amplitudo
$\frac{2F}{\abs{\omega^2-\gamma^2}}$ --- yang besar, tetapi terbatas.

\textbf{15.} Untuk $x_p = \frac{F}{2\omega}t\sin\omega t$:
\[
x_p'' = \frac{F}{2\omega}\bigl(2\omega\cos\omega t -
\omega^2t\sin\omega t\bigr)
= F\cos\omega t - \omega^2x_p :
\]
jadi sebuah penyelesaian. Lalu pada $t$ yang tetap, dengan melewatkan $\gamma \to \omega$ pada
pertanyaan 14:
\[
\frac{2F\sin\frac{(\omega-\gamma)t}2
\sin\frac{(\omega+\gamma)t}{2}}
{(\omega-\gamma)(\omega+\gamma)}
\longrightarrow
\frac{2F\cdot\frac{(\omega-\gamma)t}2\big/(\omega-\gamma)
\cdot\sin\omega t}{2\omega}
= \frac{F\,t\sin\omega t}{2\omega} .
\]
Amplitudonya tumbuh linear tanpa batas: itulah malapetaka
\omterm{pb:b2:diffeq:1}{resonansinya} --- sebab itulah para prajurit memutus langkah serentak di jembatan.

\textbf{16.} Gelombang perseginya mengusung harmonik pada setiap
frekuensi ganjil $n = 1, 3, 5, \dots$; jadi berkat kelinearannya, tiap harmonik
$n$ diperkuat tanggapan osilatornya di $\gamma = n$.
Untuk $\omega = 3$ harmonik ketiganya tepat mengenai \omterm{pb:b2:diffeq:1}{resonansinya}.
Adapun para insinyur takut pada masukan persegi (dan gigi gergaji) karena keduanya
membangkitkan \emph{semua} harmonik ganjilnya sekaligus: jadi berapa pun frekuensi
alami strukturnya, ada harmonik yang menantikannya.

\textbf{17.} Berlaku $W' = y_1y_2'' - y_1''y_2 = -qy_1y_2 + qy_1y_2 =
0$: jadi $W$ bersifat tetap (yakni Liouville dengan matriks pendamping
bertrace nol). Lalu $W = 0$ di satu titik membuat data awal $y_2$
sebanding dengan data awal $y_1$, sehingga $y_2$ sebanding dengan
$y_1$ (berkat ketunggalannya); jadi $W \neq 0$ bila dan hanya bila keduanya bebas. Bila $y(t_0) =
y'(t_0) = 0$ maka $y \equiv 0$ (berkat ketunggalannya): jadi penyelesaian tak
nol punya nol yang sederhana, dan nol yang sederhana bersifat terpencil
(sebab $y'$ bertanda tetap di dekatnya).

\textbf{18.} Di antara nol berurutan $a < b$, fungsi $y_1$ menjaga
satu tanda, katakanlah $y_1 > 0$ pada $\intoo ab$: maka $y_1'(a) > 0$ dan
$y_1'(b) < 0$ (berkat nolnya yang sederhana). Menilai konstanta $W =
y_1y_2' - y_1'y_2$ di $a$ dan $b$ memberi
\[
W = -y_1'(a)\,y_2(a) = -y_1'(b)\,y_2(b) ,
\]
jadi $y_2(a)$ dan $y_2(b)$ bertanda berlawanan (sebab $W \neq 0$
melarang keduanya lenyap): sehingga $y_2$ lenyap di $\intoo ab$
(lewat nilai antaranya). Dan ia tak dapat lenyap dua kali di sana: sebab dua nol
milik $y_2$ akan mengapit sebuah nol milik $y_1$ lewat hujah yang sama dengan
peran yang dipertukarkan, yang bertentangan dengan keberurutannya: jadi tepat
satu nol --- itulah penyisipannya.

\textbf{19.} Andaikan $z$ tak punya nol di $\intoo ab$; lalu dengan mengganti
$y, z$ dengan negatifnya, andaikan $y > 0$ dan $z > 0$ pada
$\intoo ab$. Tetapkan $\varphi = yz' - y'z$: maka $\varphi' = yz'' -
y''z = (q_1 - q_2)\,yz \leq 0$ pada $\intoo ab$, jadi $\varphi$ bersifat
tak naik. Tetapi $\varphi(a) = -y'(a)z(a) \leq 0$ (sebab
$y'(a) > 0$ dan $z(a) \geq 0$) dan $\varphi(b) = -y'(b)z(b) \geq
0$ (sebab $y'(b) < 0$ dan $z(b) \geq 0$): jadi fungsi tak naik yang
berjalan dari $\leq 0$ ke $\geq 0$ lenyap secara identik, sehingga
$(q_1 - q_2)yz \equiv 0$ pada $\intoo ab$. Bila $q_1 < q_2$
di suatu tempat di $\intoo ab$, ini mustahil (sebab $y, z > 0$ di sana):
jadi $z$ mesti lenyap tegas di dalamnya. Secara umum (dengan $q_1 \leq q_2$),
entah $z$ lenyap di $\intoo ab$, atau $\varphi \equiv 0$
memaksa $z$ sebanding dengan $y$, yang lenyap di $a$ dan $b$:
jadi pada segala kasusnya $z$ punya nol di $\intcc ab$.

\textbf{20.} Batas atasnya: bandingkanlah $y$ (dengan koefisien $q \geq
m^2$) dengan $u(t) = \sin(m(t - a))$ (dengan koefisien $m^2 \leq q$,
jadi $y$ memainkan peran $z$ pada pertanyaan 19): bila $y$ tak punya
nol di $\intoc{a}{a + \pi/m}$, maka nol $a$ dan $a +
\frac\pi m$ milik $u$ akan berurutan dengan $y \neq 0$
di antaranya, yang bertentangan dengan pertanyaan 19: jadi nol berurutan
milik $y$ berjarak $\leq \frac\pi m$. Batas bawahnya: bila dua
nol berurutan $a < b$ milik $y$ punya $b - a < \frac\pi M$,
maka $z(t) = \sin(M(t-a))$ (dengan koefisien $M^2 \geq q$) mestilah
lenyap di $\intcc ab \subset \intoo{a}{a + \pi/M}
\cup\{a\}$, tempat satu-satunya nolnya adalah $a$ itu sendiri --- tetapi
pertanyaan 19 yang diterapkan pada $\intoo{a}{b}$ dengan ketegasan di
ujungnya memberi sebuah nol di $\intcc ab$, sedangkan $z > 0$ pada
$\intoc ab$: jadi kontradiksi. Karena itu $\frac\pi M \leq b - a \leq
\frac\pi m$; dan untuk $q \equiv \omega^2$ kedua batasnya runtuh menjadi
jarak eksak $\frac\pi\omega$ milik osilator
harmoniknya.

\textbf{21.} Dengan $u = ty$: berlaku $u'' = ty'' + 2y'$, jadi $ty'' + 2y'
+ ty = u'' + u = 0$: sehingga $u = A\sin t + B\cos t = R\sin(t +
\varphi)$, lalu $y = \frac{u}{t}$ memulihkan $\frac{\sin t}t$ dan
$\frac{\cos t}t$ (\cref{exo:b2:diffeq:7}) tanpa penurunan
orde. Adapun nol setiap penyelesaian tak nolnya adalah nol
$R\sin(t + \varphi)$: yang berjarak tepat $\pi$ --- itulah
filsafat Sturm yang beraksi: bahwa koefisien $q \equiv 1$ mendiktekan
nolnya, ada rumus atau tidak.

\textbf{22.} Tetapkan $x(t) =
\frac1\omega\int_0^t\sin(\omega(t-s))F(s)\dd s$. Maka $x(0) =
0$;
\[
x'(t) = \frac1\omega\sin(0)F(t) +
\int_0^t\cos(\omega(t-s))F(s)\dd s
= \int_0^t\cos(\omega(t-s))F(s)\dd s ,
\]
jadi $x'(0) = 0$; dan $x''(t) = F(t) -
\omega\int_0^t\sin(\omega(t-s))F(s)\dd s = F(t) - \omega^2x(t)$
(lewat penurunan \omterm{thm:b2:integration:continuity}{integral berparameter} dengan batas yang berubah,
seperti pada bab integrasi). Lalu dengan $F(s) = F\cos\omega s$,
hasil-kali-ke-jumlah memberi
\[
\int_0^t\sin(\omega(t-s))\cos(\omega s)\dd s
= \frac12\int_0^t\bigl(\sin\omega t + \sin(\omega t -
2\omega s)\bigr)\dd s
= \frac{t}{2}\sin\omega t ,
\]
(sebab keping keduanya berintegral nol), jadi $x =
\frac{F}{2\omega}t\sin\omega t$: yakni pertanyaan 15 lagi, kini lewat
Duhamel.

\textbf{23.} Dalam bentuk sistemnya $X' = AX + (0, F(t))^{\mathsf T}$
dengan $\operatorname{Sp}A$ berbagian real negatif (yakni $\zeta >
0$): variasi konstanta dan \cref{exo:b2:diffeq:8}
(yakni $\vertiii{\eu^{tA}} \leq C\eu^{-\alpha t}$) memberi
\[
\norm{X(t)} \leq C\eu^{-\alpha t}\norm{X_0}
+ \int_0^t C\eu^{-\alpha(t-s)}\norm{F}_\infty\dd s
\leq C\norm{X_0} + \frac{C\norm F_\infty}{\alpha} :
\]
jadi masukan terbatas, keluaran terbatas --- secara seragam dalam data
awalnya sesudah transiennya.

\textbf{24.} Untuk $\zeta = 0$ dan $\gamma \neq \omega$,
penyelesaian pertanyaan 14 bersifat terbatas, dan menambahkan sembarang penyelesaian
homogen (yang terbatas, sebab orbit pusatnya berupa lingkaran) menjaganya tetap
terbatas; sedangkan di $\gamma = \omega$, pertanyaan 15 tumbuh linear. Jadi
bagi osilator tak teredamnya, keterbatasan di bawah pemaksaan
berkala gagal persis pada satu frekuensi --- yakni \omterm{pb:b2:diffeq:1}{resonansinya} ---
sedangkan pertanyaan 23 menunjukkan bahwa peredaman positif yang mana pun memulihkan
keterbatasannya bagi \emph{setiap} masukan yang terbatas.

\textbf{25.} (i) Bidang trace--determinan menggolongkan semua
dinamika linear bidang yang otonom, dan osilator Bagian II
menyusuri satu garis tegaknya; sedangkan pemaksaannya meninggalkan bidangnya
(sebab menjadi tak otonom), lalu Duhamel yang mengambil alih. (ii) Adapun $\gamma_*$ adalah
frekuensi yang disukai sistemnya, $A(\gamma_*)$ harga
membangkitkannya, dan $\frac{1}{2\zeta}$ faktor penguatannya
--- yakni ketajaman \omterm{pb:b2:diffeq:1}{resonansi} yang disebut para insinyur sebagai faktor
mutu. (iii) Teorema Sturm membaca ayunan dari tanda
dan besar $q$ belaka: jadi keduanya mengatur persamaan (Bessel,
Schrödinger) yang penyelesaiannya tak punya rumus dasar.
(iv) Adapun \omterm{def:b2:diffeq:wronskian}{wronskian} yang tetap (pertanyaan 17, lewat Liouville) dan
kestabilan masukan-terbatas (pertanyaan 23) dipakai ulang secara diam-diam
setiap kali buku ini bertemu persamaan berkoefisien berubah atau
sistem yang terusik.

\end{solution}
